(1) 在正入射无磁介质中光学导纳与折射率 n 成正比；用规范化导纳 \(\eta =n\) 表示相干传输的电、磁场幅。光学相位 \(\delta =2\pi nd/\lambda \)，单层的特征矩阵为
\[
M(n,\delta)=\begin{pmatrix}
\cos\delta&i\sin\delta/n\\ i n\sin\delta&\cos\delta
\end{pmatrix}.
\]
设计波长处 \(nd=\lambda _{0}/4\)，故 \(\delta =\pi /2\)，矩阵变成 \(M(n)=\begin{pmatrix}0&i/n\\in&0\end{pmatrix}\)。

(2) 对入射侧先 H 后 L 的一对层，
\[
M_{HL}=M(n_H)M(n_L)
=\begin{pmatrix}-n_L/n_H&0\\0&-n_H/n_L\end{pmatrix}.
\]
N 对后的矩阵为 \(\mathrm{diag}[(-n_L/n_H)^N,(-n_H/n_L)^N]\)。在基底导纳 ns 终止的入射端有效导纳由 \(Y_{\rm in}=(C+Dns)/(A+Bns)\) 给出，其中 A,B,C,D 为总特征矩阵元素。因此
\[
Y_{\rm in}=n_s\left(\frac{n_H}{n_L}\right)^{2N}.
\]

(3) 入射介质导纳 \(n_{0}\)，电场反射幅度为
\[
r_N=\frac{n_0-Y_{\rm in}}{n_0+Y_{\rm in}},\qquad
R_N=|r_N|^2=
\left[\frac{n_0-n_s(n_H/n_L)^{2N}}
{n_0+n_s(n_H/n_L)^{2N}}\right]^2.
\]
当 nH>nL 且 N 越大，有效导纳增长，\(r_N\to-1,R_N\to1\)，理想多层结构形成高反射带；反射增强来自多界面相位相干，而不是材料吸收。

(4) 若 N=0，没有任何薄膜，\(Y_{\rm in}=n_s\)，恢复单界面反射率 \([(n_0-n_s)/(n_0+n_s)]^2\)。偏离设计波长时每层 \(\delta \) 不再是 \(\pi /2\)，矩阵将出现非零对角余弦项和频率依赖复耦合，反射率随波长起伏形成有限宽度的 Bragg 停带，不能认为上述设计波长公式对所有颜色都严格成立。
